\documentclass[11pt]{article}
\usepackage{amsmath,amssymb,amsthm,hyperref}
\newtheorem{proposition}{Proposition}
\newcommand{\F}{\mathbb F}
\title{Pairwise-valid context mechanisms need not compose:\\a finite braid obstruction and minimal action repair}
\author{Contextual composition programme}
\date{13 November 2026}
\begin{document}\maketitle
\begin{abstract}
We give a small exact obstruction to composing independently supplied context
mechanisms. Each of two mode mechanisms crosses a resource accumulator by a
valid distributive law, and the mode mechanisms themselves have the usual
independent swap. Nevertheless the three crossings fail their braid equation;
the proposed combined multiplication is not associative. In the displayed
two-bit model, replacing independent mode parities by their generated action
group repairs composition, with six mode states and 24 total states. These
counts are minimal for the specified observations. This is a worked strict
writer/accumulator specialization of classical distributive-law mathematics,
not a new general composition theorem for arbitrary contextads.
\end{abstract}

\section{What is being combined?}
The catalogue's composed-contextads direction asks for coherent combinations
of mechanisms, not just several pieces of data in a tuple. The existing
\emph{One relative construction for context, content, and effects}, section
``One mixed multiplication from a coherent crossing'', explicitly assumes a
complete supplied pseudodistributive law. It does not assert that arbitrary
independently valid mechanisms admit one. Here we isolate an additional
problem when there are three mechanisms: pairwise laws are insufficient.

Our strict setting uses accumulator grades. A monoid $M$ defines the writer
monad $W_M X=M\times X$ on sets, with unit $(1,x)$ and multiplication
$(m,(n,x))\mapsto(mn,x)$. A crossing between grade monoids satisfying the two
unit and two multiplication equations extends naturally to a distributive
law of their writer monads. The payload $x$ is untouched. This supplies an
ordinary compositional semantics with explicit accumulated grades; the
resulting grade monoid can also be regarded as a one-object category. We do
\emph{not} identify every writer monad with the programme's general native
contextad presentations, nor supply dependent reindexing, nontrivial grade
cells, or their pseudocoherences here.

The familiar higher compatibility condition is the Yang--Baxter equation
for iterated distributive laws, treated generally by Cheng
\cite{Cheng}. Our contribution here is the explicit failure, repair, and
observable-state accounting for this finite application; no novelty is
claimed for the underlying theory.

\section{An exact compatibility criterion}
Let $A$ be a group, written multiplicatively for this section, and let $P,Q$
be involutive automorphisms of $A$. Let $B=C=\mathbb Z/2$ be mode grades.
Define pairwise crossings toward normal order $ABC$ by
\[
 BA\longrightarrow AB:(b,a)\longmapsto(P^b(a),b),\qquad
 CA\longrightarrow AC:(c,a)\longmapsto(Q^c(a),c),
\]
and $CB\to BC:(c,b)\mapsto(b,c)$. The last crossing says that the two
mode mechanisms are independent.

\begin{proposition}[Pairwise validity]
All three crossings are distributive laws of the corresponding writer
monads, regardless of whether $P$ and $Q$ commute.
\end{proposition}
\begin{proof}
For $BA\to AB$, the unit equations are $P^0a=a$ and $P^b1=1$.
The two multiplication equations are
\[
 P^b(ad)=P^b(a)P^b(d),\qquad
 P^{b+e}(a)=P^b(P^e(a)),
\]
where mode addition is modulo two. They follow from automorphism and
involution. The $Q$ case is identical, and swap is the product-monoid law.
Naturality in the untouched payload is immediate.
\end{proof}

\begin{proposition}[Three-way obstruction]
The three crossings satisfy the braid equation if and only if $PQ=QP$.
Equivalently, the attempted multiplication on $A\times B\times C$,
\begin{equation}\label{eq:attempt}
 (a,b,c)(d,e,f)=(aP^bQ^c(d),\ b+e,\ c+f),
\end{equation}
is associative if and only if $P$ and $Q$ commute.
\end{proposition}
\begin{proof}
The two rewriting paths from $CBA$ to $ABC$ send $(c,b,a)$ to
$(Q^cP^b(a),b,c)$ and $(P^bQ^c(a),b,c)$. They agree exactly when
$P$ and $Q$ commute. If they do, $(b,c)\mapsto P^bQ^c$ is an action
of $B\times C$ by automorphisms, and \eqref{eq:attempt} is its semidirect
product. For necessity, use $x=(1,0,1)$, $y=(1,1,0)$, and $z=(d,0,0)$.
The $A$ coordinates of $(xy)z$ and $x(yz)$ are $PQ(d)$ and $QP(d)$.
\end{proof}

Thus the defect is not an implementation race: the proposed algebra fails
associativity before scheduling or probabilities enter. A compilation rule
that moves mode instructions past accumulated updates must check this
three-way condition, even after all pairwise tests pass.

\section{Two binary wires: a concrete failure}
Take $A=V=\F_2^2$, with vector addition. Translations by $a\in V$
toggle a specified pair of binary wires. Mode $P$ swaps the wire roles;
mode $Q$ performs the reversible linear recoding that XORs the second wire
into the first:
\[
 P(x,y)=(y,x),\qquad Q(x,y)=(x+y,y).
\]
Both square to the identity and preserve vector addition. Moving a linear
recoding past a toggle changes the toggle mask by that recoding, precisely
the crossing above. Each isolated mechanism therefore has coherent
accumulation semantics. But they are not independent modes:
\[
 PQ(1,0)=(0,1),\qquad QP(1,0)=(1,1).
\]
This difference is observable on actual wire values, not just in the choice
of notation. If $x=(0,0,1)$, $y=(0,1,0)$, and
$z=((1,0),0,0)$, interpreting the first $0$ as the zero vector, then
\eqref{eq:attempt} assigns different toggle masks to its two bracketings.
Among the 4096 triples in the attempted 16-element structure, the exact
test finds 768 associativity failures.

There is a second intuitive diagnosis. Recording only how many $P$ and $Q$
instructions occurred modulo two identifies the empty history with $PQPQ$.
Yet $(PQ)^2\ne1$. Reordering and accumulating the separate mode parities
has thrown away a cross-mechanism effect. Keeping the chronological mode
word would work but is unnecessarily large for the present observations.

\section{A repair that retains exactly the observable action}
Let $G=\langle P,Q\rangle\leq\operatorname{Aut}(V)$ be the group of actual
linear transformations generated by these instructions. The three nonzero
vectors of $V$ are permuted by $G$. Here $P$ exchanges $(1,0)$ and $(0,1)$,
while $Q$ exchanges $(0,1)$ and $(1,1)$. These transpositions generate every
permutation of the three nonzero vectors, so $G\cong S_3$ and $|G|=6$.
The generated context is
\[
 H=V\rtimes G,\qquad (a,g)(d,h)=(a+g(d),gh).
\]
Products of functions are composed right to left. It acts on wire values by
$F_{a,g}(v)=a+g(v)$, with $F_{(a,g)(d,h)}=F_{a,g}\circ F_{d,h}$.
Thus our convention executes the right factor first. Reversing the temporal
convention consistently gives the same obstruction and repair.

\begin{proposition}[Repair and minimality at the stated observation scope]
$H$ is an associative 24-element context monoid. Any exact summary of mode
histories from which one can recover their action on every $v\in V$ has
at least six states. Any exact summary of combined toggle and mode histories
recovering their action on every $v\in V$ has at least 24 states. Both bounds
are attained by $G$ and $H$, respectively.
\end{proposition}
\begin{proof}
The semidirect formula is associative because $g$ is additive and
$(gh)(v)=g(h(v))$; $(0,1)$ is the unit. There are six distinct generated
mode functions, so any summary identifying two of them loses some probe
value. Every $(a,g)$ is generated by toggles and modes. If
$F_{a,g}=F_{d,h}$ on all probes, evaluating at zero gives $a=d$, and then
subtracting $a$ gives $g=h$. Hence all 24 affine behaviors are distinct,
giving the lower bound. Their multiplication supplies attaining summaries.
\end{proof}

Minimality means distinguishability under arbitrary initial wire probes
and exact output observation, not minimum memory under a fixed input,
restricted tests, approximation, or a demand to retain the entire history.
For example, observing zero only cannot distinguish linear modes. The
repair deliberately forgets different mode words with the same total action;
that is sound for these observations but would not preserve provenance,
instruction costs, or intermediate traces. Including those observations
would require a different summary.

The general repair pattern is to retain the image of the mode-word monoid
in $\operatorname{Aut}(A)$ and form the corresponding semidirect product.
It is not always finite. Here it is finite, and preserves the meaning of
each individual mechanism but rejects their asserted independence. It does
not magically provide a law combining the original independent mode bits.

\section{Checks and limitations}
\texttt{check\_composed.py} checks both pairwise multiplication equations,
the unit equations, all 16 braid inputs, all 4096 attempted-product
associativity cases, all 13824 repaired-product associativity cases, and the
24 distinct affine observations. As a positive control it replaces $Q$ by
the identity and checks all 4096 associative triples in the resulting
independent-mode construction. The output is in
\texttt{check-results.json}; no external Python dependencies are needed.

The written propositions give the general criterion; finite enumeration is
regression support, not their substitute. No Lean formalization or new
literature/novelty claim is made. In particular, this example does not solve
the general composed-contextad direction, construct a pseudodistributive
law in the native dependent setting, or show that all possible combinations
of these mechanisms are impossible. It rules out precisely the simultaneous
specified pairwise crossings and independent swap. The positive repair
shows what additional interaction data make this instance compositional.

\begin{thebibliography}{9}
\bibitem{Cheng} Eugenia Cheng, \emph{Iterated distributive laws},
Mathematical Proceedings of the Cambridge Philosophical Society 150 (2011),
459--487. \url{https://arxiv.org/abs/0710.1120}.
\end{thebibliography}
\end{document}
